% Generated draft: requires mathematical and source review.
\documentclass[11pt]{article}
\usepackage[margin=25mm]{geometry}
\usepackage{fontspec}
\setmainfont{Latin Modern Roman}
\usepackage{amsmath,amssymb,mathtools,unicode-math}
\setmathfont{Latin Modern Math}
\usepackage{enumitem,needspace,xurl,hyperref}
\setlength{\parindent}{0pt}
\setlength{\parskip}{5pt}
\setlength{\emergencystretch}{3em}
\newcommand{\N}{\mathbb{N}}
\newcommand{\Z}{\mathbb{Z}}
\newcommand{\Q}{\mathbb{Q}}
\newcommand{\R}{\mathbb{R}}
\newcommand{\C}{\mathbb{C}}
\newcommand{\F}{\mathbb{F}}
\newcommand{\A}{\mathbb{A}}
\newcommand{\PP}{\mathbb{P}}
\newcommand{\E}{\mathbb{E}}
\newcommand{\eps}{\varepsilon}
\newcommand{\dd}{\,\mathrm{d}}
\begin{document}
\section*{UnsolvedMath ID: 30005591}
\textbf{Source problem code:} catalogue entry
\subsection{Definitions and mathematical statement}
A steady gradient Ricci soliton is a Riemannian manifold $(M,g)$ equipped with a smooth function $f$ satisfying
\[
\operatorname{Ric}_g+\nabla_g^2 f=0.
\]
Here the sign of $f$ is a convention. Two such metrics $(M_1,g_1)$ and $(M_2,g_2)$ are isometric up to rescaling if there exist a constant $c>0$ and a diffeomorphism $\Phi:M_1\to M_2$ such that $\Phi^*g_2=cg_1$. The proposed conclusion concerns the metrics; the source does not require an identification of the soliton potentials.

Let $\mathcal F_3$ denote the class called three-dimensional flying-wing steady gradient Ricci solitons in the supplied source, and let $\theta(S)$ denote the asymptotic cone angle of a member $S$. The principal conjecture is
\[
\forall S_1,S_2\in\mathcal F_3,\qquad
\theta(S_1)=\theta(S_2)
\ \Longrightarrow\ 
\exists c>0\ \exists\Phi:M_1\xrightarrow{\cong}M_2,
\quad \Phi^*g_2=cg_1.
\]
This is a conditional formalization using the source terminology: the supplied record does not define the flying-wing class, the asymptotic cone construction, the angle convention, or its admissible range. Those definitions must be recovered from the cited sources before this becomes a fully specified conjecture. No additional completeness, topology, curvature, or noncollapsing hypothesis is silently imposed here.

For the accompanying symmetry questions, an ancient Ricci flow is a smooth family of Riemannian metrics satisfying
\[
\partial_t g(t)=-2\operatorname{Ric}_{g(t)}
\]
on a time interval unbounded below. The orthogonal group is $O(k)=\{A\in\operatorname{Mat}_{k\times k}(\R):A^{\mathsf T}A=I_k\}$. An isometric action on $(M,g)$ is a smooth group action whose elements preserve $g$. Positive curvature operator means that, at every point $p$,
\[
\langle\mathcal R_p\omega,\omega\rangle>0
\qquad\text{for every }0\ne\omega\in\Lambda^2T_pM,
\]
where $\mathcal R_p$ is the curvature operator with the convention that the round sphere has positive curvature operator.

The source additionally asks:
\begin{enumerate}
\item Does $O(2)$-symmetry hold for three-dimensional ancient Ricci flows?
\item Does $O(2)$-symmetry hold for steady gradient Ricci solitons of dimension greater than three with positive curvature operator?
\item Does $O(3)$-symmetry hold for all noncollapsed four-dimensional steady gradient Ricci solitons?
\end{enumerate}
These are source questions with unresolved formulation requirements. The intended rotational actions, their effectivity and orbit structure, and whether a single action must preserve an ancient flow at every time are unspecified. Allowing arbitrary actions would permit the trivial action and would not express the intended symmetry claim. The meaning and scale range of noncollapsed are also unspecified. Completeness and any additional curvature, topology, boundedness, or noncollapsing assumptions for the first two questions are not supplied. The four-dimensional question does not explicitly repeat the positive-curvature-operator assumption, so that assumption has not been transferred to it.
\subsection{Short English statement}
The main conjecture asks whether a three-dimensional flying-wing steady gradient Ricci soliton is determined, up to isometry and constant metric rescaling, by its asymptotic cone angle. Three further questions concern orthogonal symmetry of ancient Ricci flows and steady gradient Ricci solitons. The definitions of flying wings, cone angle, noncollapsing, and the intended symmetry actions require clarification from the sources. This is a statement-only draft requiring human mathematical review.
\subsection{Sources}
\begin{itemize}
\item Supplied local UnsolvedMath registry record, permanent ID $30005591$, title \emph{Rigidity and Symmetry of Steady Ricci Solitons}, exact source problem code \texttt{catalogue entry}, in the Oberwolfach Reports Open Problems collection. The record lists proposed year 2023 and supplies no named proposer.
\item Supplied report citation: \emph{Partial Differential Equations (2024)}, Oberwolfach Reports, pp.~1789--1842, DOI 10.4171/owr/2023/32. Source URL and literature-source label \emph{Oberwolfach Report 32}: \url{https://doi.org/10.4171/owr/2023/32}. The citation's 2024 year and the DOI's 2023 component are retained as supplied; their relationship has not been verified.
\item Supplied literature source labeled \emph{Lai primary article}: \url{https://arxiv.org/abs/2207.02714}. The supplied assessment attributes to Lai existence of three-dimensional flying wings for every cone angle and their symmetry, and describes uniqueness for a fixed angle modulo rescaling as conjectural. The record does not specify the admissible angle range or the precise hypotheses of these attributed results.
\item Literature assessment date: \textbf{2026-08-22}. The supplied assessment labels the problem open and says its literature triage was derived from the supplied Oberwolfach curation. This is supplied dated curation, not an independently verified assessment of the literature or current status.
\item The registry labels the statement \texttt{corrected\_verified} and says it was reconstructed and checked against the official Oberwolfach report. This is supplied verification metadata, not verification performed for this draft.
\item The extracted original text also attributes a classification of noncollapsed translators in $\R^4$ to Choi--Haslhofer--Hershkovits, citing local reference labels [3, 2], and mentions a classification of convex translators in $\R^3$ under reference label [4]. Bibliographic entries for those labels are not supplied. These are contextual source attributions, not independently verified claims.
\end{itemize}
\end{document}
